\documentclass[10pt,a4paper]{amsart}
\usepackage[margin=19mm]{geometry}
\usepackage{amsmath,amssymb,mathtools}
\usepackage[scr=rsfs,cal=euler]{mathalfa}
\usepackage{xfrac}
\usepackage[unicode,hidelinks,bookmarks=false]{hyperref}
\newcommand{\ie}{i.e.\ }
\newcommand{\cf}{cf.\ }
\newcommand{\resp}{respectively\ }
\newcommand{\Subsection}{\subsection}
\providecommand{\nobreakdash}{\nobreak\hbox{-}}
\setlength{\emergencystretch}{2em}
\multlinegap0pt
\usepackage{bm}
\usepackage{bbm}
\usepackage{dsfont}
\usepackage{xspace}
\usepackage{cancel}
\usepackage{mathtools}
\usepackage{amsthm}
\makeatletter
\@ifundefined{theorem}{\newtheorem{theorem}{Theorem}}{}
\@ifundefined{lemma}{\newtheorem{lemma}[theorem]{Lemma}}{}
\@ifundefined{proposition}{\newtheorem{proposition}[theorem]{Proposition}}{}
\makeatother
\newcommand{\R}{\mathbb{R}}
\title[Scattering Theory For 3D Cubic Damped Magnetic Schr\"odinger]{Scattering Theory For 3D Cubic Damped Magnetic Schr\"odinger Equation}
\author{Luca Fanelli, Haruya Mizutani, Yilin Song, Ying Wang,, Jiqiang Zheng}
\date{Source: arXiv:2608.11834v1 (CC BY 4.0)}
\begin{document}
\maketitle

\noindent\emph{Source and licence.} Scattering Theory For 3D Cubic Damped Magnetic Schr\"odinger Equation. Version arXiv:2608.11834v1 (CC BY 4.0), distributed under the Creative Commons Attribution 4.0 International licence, as recorded in the arXiv licence metadata for this exact version and shown on the abstract page https://arxiv.org/abs/2608.11834v1. Full rights notice: licenses/arxiv-2608.11834-CC-BY-4.0.txt.\par\smallskip

\noindent\textit{Editorial scope.} Counted unit: the Proposition whose source label is \texttt{prop:uniformenergy} in the article (source0.tex lines 1198--1207, source label \texttt{prop:uniformenergy}), delivered together with the article's own complete original proof of it (source0.tex lines 1210--1294, one \texttt{proof} environment of 85 lines). No mathematics is written, repaired or re-derived: statement and proof are the article's own text line for line. Required context retained uncounted: equ:DMNLS (lines 134--139); equ:Gpd (lines 143--145); lemma\_2\_2 (lines 521--526); prop:LWP (lines 598--608); prop:CGT (lines 669--676); equ:uniformenergy\_2 (lines 672--675); proof\_energy\_7 (lines 1171--1174). Every definition, display and label of those lines is retained unchanged. Omitted from this package and not redistributed: 1--133, 140--142, 146--520, 527--597, 609--668, 676--1170, 1175--1197, 1208--1209, 1295--2777. Nothing inside a retained excerpt is omitted or altered: every retained body is delivered exactly as the article's lines are written, so no omission receipt of any kind accompanies this package. Citations: the retained excerpts carry no citation command at all, so no bibliography entry of the article is transcribed and no reference is redistributed. Reference adaptation: none was needed and none was made; every reference inside the delivered unit resolves inside it, so no unresolved reference or citation is produced. Printed slips kept literal: no printed word, symbol, hypothesis or number of the counted statement or of its proof was altered. Layout and class: the article's own class is \texttt{amsart} with packages including graphicx, amssymb, amsmath, amsthm, natbib, color, amscd, hyperref, geometry, comment, tikz; the wrapper is the lane's pinned \texttt{amsart} editorial wrapper at 10pt on A4, which restates the theorem-like environments the retained text needs and adds the article's own macro definitions that the retained text needs (\texttt{\textbackslash R}), with extra wrapper packages (bm, bbm, dsfont, xspace, cancel, mathtools). The delivered numbering is the unit's own. Assets: no retained span contains a figure environment, a graphics-inclusion command, a PDF-page inclusion, a table or a TikZ picture; no image file is copied into this package, the package declares no asset dependency and no image file is redistributed with it. Licence: version arXiv:2608.11834v1 of this article is distributed under the Creative Commons Attribution 4.0 International licence, as recorded in the arXiv licence metadata for this exact version and shown on the abstract page https://arxiv.org/abs/2608.11834v1. A full rights notice is delivered at licenses/arxiv-2608.11834-CC-BY-4.0.txt. Characters: every retained excerpt is pure ASCII, so the delivered engine, XeLaTeX with its default Latin Modern text font, reports no missing character.
	\begin{equation}\label{equ:DMNLS}
		\begin{cases}
			i\partial_tu+\Delta_{G}u+iau=|u|^{2}u,\quad (t,x)\in\R^+\times\R^3,\\
			u(0,x)=u_0(x),
		\end{cases}
	\end{equation}

	\begin{equation}\label{equ:Gpd}
		C|\xi|^2\geq G(x)\xi\cdot\xi\geq c|\xi|^2. 
	\end{equation}

\begin{lemma}
	\label{lemma_2_2}
	For any $s\in \R$ and any $1<p<\infty$, the spaces $W^{s,p}$, $W_A^{s,p}$, and $W_G^{s,p}$ coincide with equivalent norms.
	%		For $s\in[0,2]$, we have the following equivalence of Sobolev norms:		\begin{align*}			\|f\|_{{H}^s_{G}}\simeq\|f\|_{{H}_{A}^s}\simeq\|f\|_{H^s}. 		\end{align*}
	%		In particular, we have		\begin{align*}			\|f\|_{{H}_{G}^1}\simeq\| f\|_{{H}_A^1},\,\,\|\langle\Delta_{G}\rangle f\|_{L^2}\simeq\|\langle \Delta_A\rangle f\|_{L^2}.		\end{align*}
\end{lemma}

\begin{proposition}[Local well-posedness]\label{prop:LWP}
	Let $0<\epsilon<\frac12$, $2< q<\frac2{1-2\epsilon}$ and $u_0\in H^{1+\epsilon}(\R^3)$. Then there exist $T>0$ and a unique solution $u\in C([0,T], H^{1+\epsilon}(\R^3))\cap L^q([0,T], L^\infty(\R^3))$ to \eqref{equ:DMNLS} satisfying the following statements: 
	\begin{itemize}
		\item The map $u_0\in H^{1+\epsilon}(\R^3)\mapsto u\in C([0,T], H^{1+\epsilon}(\R^3))$ is continuous. 
		\item If we denote by $T_{\rm max}$ the maximal forward time of existence, then  either  $T_{\rm max}=+\infty$ or $T_{\rm max}<+\infty$ and
		\begin{equation*}%\label{equ:blowcri}
			\lim_{t\to T_{\rm max}}\|u(t)\|_{H^{1+\epsilon}(\R^3)}=+\infty.
		\end{equation*}
		\item If $u_0\in H^r(\R^3)$ for some $r>1+\epsilon$, then $u\in C([0,T],H^r(\R^3))$. 
	\end{itemize}
\end{proposition}

\begin{proposition}[Global existence]\label{prop:CGT}
	%Let $0<\epsilon<1$. 			Assume that the following holds: there exists a continuous function $\Psi:\R_+\to\R_+^\ast$ \begin{color}{blue}[Mizutani: (1) Using the symbol $C$ would confuse readers. (2) Please write explicitly the definition of $\R_+^*$ instead of using the notation $\R_+^*$]\end{color} such that for any $u_0\in H^{1+\epsilon}(\R^3)$ and $T>0$, if \begin{color}{red}$u\in C([0,T],H^{1+\epsilon}(\R^3))$\end{color} is a solution to \eqref{equ:DMNLS} on $[0,T]$ with initial data $u_0$, then\begin{align}\label{prop:CGT_1}\sup_{[0,T]}\|u(t)\|_{H^1(\R^3)}\leq \Psi\big(\|u_0\|_{H^1(\R^3)}\big).\end{align}		Then, the maximal solutions to \eqref{equ:DMNLS} are global forward in time.
	For any $u_0\in H^{1+\epsilon}(\R^3)$, the maximal solution $u$ constructed in Proposition \ref{prop:LWP} exists globally forward in time, namely $u\in C([0,\infty),H^{1+\epsilon}(\R^3))$. Moreover, $u$ satisfies
	\begin{align}
		\label{equ:uniformenergy_2}
		\sup_{t\ge0}M(u(t))\le M(u_0),\quad \sup_{t\ge0}E(u(t))\lesssim E(u_0)+M(u_0).
	\end{align}
\end{proposition}

	\begin{align}
		\label{proof_energy_7}
		\sup_{R>0}M_\rho(t)\ge \delta\lambda_1+(1-2\delta)\lambda_2+\delta\lambda_3-C\{M(u_0)+E(u_0)\}
	\end{align}

\begin{proposition}[Local energy decay]\label{prop:uniformenergy} Let $u\in C([0,\infty),H^{1+\epsilon}(\R^3))$ be the global solution obtained in Proposition \ref{prop:CGT}. Then
	\begin{align}
		\label{equ:locmasdec}
		&\sup_{R>0}\frac{1}{R^2}\int_0^\infty\int_{|x|=R}|u|^2d\omega\,dt+\sup_{R>0}\frac{1}{R}\int_0^\infty\int_{|x|\le R}|{\nabla_A u}|^2\,dx\,dt\\\nonumber&\quad +\int_0^\infty\int_{\R^3}\frac{|\nabla_A^\tau u|^2}{|x|}\,dx\,dt\lesssim   E(u_0)+M(u_0). 
	\end{align}
	Moreover, for any $\chi\in C_0^\infty(\R^3)$ and $0\le s<1$, 
	\begin{equation}\label{equ:locenedecHs}
		\lim_{t\to\infty}\|\chi u(t)\|_{H^s(\R^3)}=0.
	\end{equation}
\end{proposition}

\begin{proof}
	% \eqref{equ:locmasdec} follows by letting $t\to \infty$ in \eqref{proof_energy_7}. %By \eqref{equ:locmasdec}, there exists $t_n$ with $t_n\nearrow \infty$ as $n\to \infty$ such that $\chi u(t_n)\to 0$ in $L^2$ as $\to n\to\infty$. This sequential convergence, together with $M(t)\le M(t_n)$ for any $t\ge t_n$ by \eqref{equ:mass}, shows 
%     The proof of \eqref{equ:locenedecHs} is the same as that of \cite[Corollary 3.6]{LS25}. 
% \end{proof}
% \begin{proof}
Estimate \eqref{equ:locmasdec} follows by letting $t\to\infty$ in \eqref{proof_energy_7}.
For \eqref{equ:locenedecHs},
let $\chi\in C_c^\infty(\mathbb R^3)$ be real-valued and choose
$R_\chi>0$ such that
\[
\operatorname{supp}\chi\subset B(0,R_\chi).
\]
By the coarea formula and the first term in \eqref{equ:locmasdec},
\[
\begin{aligned}
\int_0^\infty\int_{B(0,R_\chi)}|u|^2\,dx\,dt
&=
\int_0^{R_\chi}
\int_0^\infty\int_{|x|=r}|u|^2\,d\omega\,dt\,dr
\\
&\le
\frac{R_\chi^3}{3}
\sup_{R>0}\frac1{R^2}
\int_0^\infty\int_{|x|=R}|u|^2\,d\omega\,dt
<\infty.
\end{aligned}
\]
The second term in \eqref{equ:locmasdec}, evaluated at $R=R_\chi$, also gives
\[
\int_0^\infty\int_{B(0,R_\chi)}
|\nabla_Au|^2\,dx\,dt<\infty.
\]
Hence
\[
\chi u\in L^2\bigl((0,\infty);H^1(\mathbb R^3)\bigr),
\]
where we have used the boundedness of $A$ and Lemma \ref{lemma_2_2}

Set
\[
f(t):=\|\chi u(t)\|_{L^2}^2.
\]
Using equation \eqref{equ:DMNLS}, integration by parts gives the localized
mass identity
\[
\frac12 f'(t)
=
-\int_{\mathbb R^3}a\chi^2|u|^2\,dx
+
\operatorname{Im}\int_{\mathbb R^3}
\overline u\,G\nabla_Au\cdot\nabla(\chi^2)\,dx.
\]
Consequently, by the boundedness of $G$ and $a$,
\[
|f'(t)|
\lesssim_\chi
\int_{B(0,R_\chi)}
\bigl(|u(t,x)|^2+|\nabla_Au(t,x)|^2\bigr)\,dx.
\]
The preceding local spacetime estimates therefore imply
\[
f,f'\in L^1(0,\infty).
\]
Thus $f(t)$ has a finite limit as $t\to\infty$. Since
$f\ge0$ and $f\in L^1(0,\infty)$, this limit must be zero:
\[
\|\chi u(t)\|_{L^2}\longrightarrow0.
\]

Finally, the global bound \eqref{equ:uniformenergy_2}, the ellipticity condition \eqref{equ:Gpd},
and Lemma \ref{lemma_2_2} imply
\[
\sup_{t\ge0}\|\chi u(t)\|_{H^1}<\infty.
\]
Therefore, for every $0\le s<1$, interpolation yields
\[
\|\chi u(t)\|_{H^s}
\lesssim
\|\chi u(t)\|_{L^2}^{1-s}
\|\chi u(t)\|_{H^1}^{s}
\longrightarrow0
\qquad (t\to\infty).
\]
This proves \eqref{equ:locenedecHs}.
\end{proof}

\end{document}
